% CECI EST UN FICHIER DE MODÈLE LATEX POUR LES ARTICLES INCLUS DANS
% *Anthology of Computers and the Humanities*. AJOUTEZ L'OPTION
% 'final' LORS DE LA CRÉATION DE LA VERSION FINALE DE L'ARTICLE. 
% NE PAS modifier la classe du document
%\documentclass[final,fra]{anthology-ch} % pour la version finale
\documentclass[final,fra]{anthology-ch}         % pour la soumission

% CHARGER LES PACKAGES LaTeX
\usepackage{booktabs}
\usepackage{graphicx}

% AJOUTEZ vos propres packages avec \usepackage{}

% TITRE DE LA SOUMISSION
% Remplacez ceci par le titre de votre soumission
\title{Annoter et détecter les citations : vers un cadre unifié entre linguistique et humanités computationnelles}

% INFORMATIONS SUR LES AUTEURS ET LES AFFILIATIONS
% Pour chaque auteur, utilisez une nouvelle commande \author, avec
% les numéros entre crochets indiquant les affiliations correspondantes
% (section suivante) et l'ORCID-ID pour chaque auteur.  
\author[1]{Agnès Saulnier}[
  orcid=0009-0002-9715-1405,
  email=asaulnier@ina.fr
]

% Il doit y avoir un appel à \affiliation pour chaque affiliation des auteurs.
% Plusieurs affiliations peuvent être attribuées à un auteur
% et une affiliation peut être attribuée à plusieurs auteurs. 
\affiliation{1}{Institut national de l’audiovisuel, Bry-sur-Marne, France}

% MOTS-CLÉS
% Fournissez un ou plusieurs mots-clés séparés par des virgules
% en utilisant les commandes suivantes
\keywords[fra]{discours rapporté, citation, corpus annotés, corpus audiovisuel, TAL}
\keywords[eng]{reported speech, quotation, annotated corpora, audiovisual corpora, NLP}

% MÉTADONNÉES DE LA PUBLICATION
% Ces champs seront remplis lors de la publication ; 
% ils peuvent être laissés par défaut lors de la soumission
\pubyear{2026}
\pubvolume{4}
\pagestart{37}
\pageend{43}
\paperorder{5}
\conferencename{Actes de la Conférence Humanistica}
\conferenceeditors{Serena Crespi, Simon Gabay, Martin Grandjean, Ariane Pinche, Marie Puren et Léa Saint-Raymond}
\doi{10.63744/UM96NSIlUNaK}
\customhead{}

\addbibresource{bibliography.bib}

%%%%%%%%%%%%%%%%%%%%%%%%%%%%%%%%%%%%%%%%%%%%%%%%%%%%%%%%%%%%%%%%%%%%%%%%%%%
% DÉBUT DU TEXTE
\begin{document}

\maketitle

\begin{abstract}
Quotation is a central but theoretically unstable object, situated at
the intersection of linguistic, narrative, media, and computational
traditions. This plurality is reflected in existing annotated corpora,
which rely on heterogeneous scopes and categories, with consequences for
automatic detection and evaluation. Based on a review of the literature
and an analysis of existing corpora, this article shows why these
divergences become problematic in computational contexts. It then
proposes a typology of quotation grounded in linguistic descriptions,
designed to make explicit the scope choices underlying annotated
corpora. The discussion shows how this typology can guide annotation
practices, support system comparison, and facilitate extensions to
audiovisual data.
\end{abstract}

\section{Introduction et problématique}

La citation constitue un objet central des sciences humaines, mais elle
ne renvoie pas, en linguistique, à une notion théorique unifiée. Elle
est décrite à partir de traditions descriptives distinctes, mobilisant
des critères syntaxiques, énonciatifs, narratifs, pragmatiques ou
juridiques, dont les périmètres ne se recouvrent que partiellement.

Dans ce contexte, les travaux en traitement automatique du langage (TAL)
ont cherché à opérationnaliser la citation à partir de divers corpus
annotés, parmi lesquels PARC \cite{pareti2012database}, FRACAS \cite{richard2024fracas}, \textit{DE-News} \cite{petersenfrey2024dataset} et \textit{QuoteBank} \cite{vaucher2021quotebank} ont joué un rôle structurant pour la détection et
l'attribution. Ces ressources reposent toutefois sur des choix de
périmètre différents, ce qui devient problématique lorsqu'on cherche à
comparer des systèmes ou à interpréter leurs performances.

L'objectif de cet article n'est ni de proposer une définition
universelle de la citation, ni d'unifier artificiellement les corpus
existants, mais de clarifier les enjeux linguistiques et computationnels
liés à son opérationnalisation. Nous proposons une typologie conçue
comme un outil de mise en cohérence, visant à expliciter les choix de
périmètre et à fournir un cadre commun pour articuler corpus, méthodes
de détection et protocoles d'évaluation, y compris en audiovisuel.

\section{État de l'art : une notion plurielle et difficile à stabiliser}

La notion de citation, ou plus largement de discours rapporté, occupe
une place centrale dans plusieurs traditions disciplinaires sans
renvoyer à un objet théorique unifié. Selon les cadres mobilisés, elle
peut être décrite comme une structure morphosyntaxique, une opération
énonciative, un procédé narratif, un acte de langage ou un objet
juridique et médiatique. Cette pluralité se reflète dans les périmètres
hétérogènes des corpus.

Les grammaires de référence du français décrivent le discours rapporté à
partir de critères morphosyntaxiques et distinguent discours direct,
discours indirect et discours indirect libre \cite{riegel1994grammaire,grevisse2016bon}. Le discours direct et le discours
indirect sont identifiés par des marqueurs formels explicites
(ponctuation, subordination, concordance des temps), tandis que le
discours indirect libre repose sur un faisceau d'indices linguistiques
et contextuels (temps verbaux, déictiques, point de vue), sans marqueur
explicite d'attribution. Dans cette approche, la citation est définie à
partir de configurations linguistiques identifiables, sans prise en
charge systématique des opérations de reformulation ou de la prise de
position du locuteur citant.

La linguistique du discours et de l'énonciation montre que rapporter un
discours ne consiste pas à le copier mais à le reformuler, que le
locuteur citant prend position vis-à-vis du discours rapporté, et qu'il
existe des formes intermédiaires entre citation et interprétation
\cite{bakhtine1977discours,authier1996ces}, fréquentes dans le discours
journalistique \cite{rosier1998discours}.

En narratologie, le discours rapporté est analysé comme un mode de
représentation des paroles et des pensées, centré sur le point de vue et
la mise en scène des voix dans le récit \cite{genette1972figures,genette1982palimpsestes,cohn1978transparent}. L'attention porte sur la frontière entre narrateur et personnage,
et sur des formes telles que le discours indirect libre ou la pensée
représentée.

En pragmatique, cette réflexion est prolongée par la distinction entre
\textit{oratio recta}, reprise d'un dire comme tel, et \textit{oratio obliqua}, reprise
d'un contenu reformulé \cite{austin1962how,searle1969speech,recanati2000oratio}.
Cette opposition permet de distinguer la reprise locale d'un dire
formulé de la circulation plus abstraite de contenus propositionnels,
indépendamment de leur énonciation d'origine.

Enfin, les cadres juridiques définissent la citation comme une reprise partielle d’un extrait préexistant, intégrée dans une œuvre seconde et justifiée par une finalité (critique, scientifique, etc.)\footnote{Code de la propriété intellectuelle, art. L.122-5.}. La doctrine du droit d’auteur souligne en outre que cette reprise doit demeurer subordonnée au propos de l'œuvre citante et ne pas se substituer à l’œuvre source \cite{vivant_droit_2009,chambathouillon2016droit}. Les études médiatiques, de leur côté, s'intéressent à la circulation de fragments discursifs autonomisés (slogans ou petites phrases) dans l’espace public \cite{kriegplanque2010formule,boyer2018petites}.

%Enfin, les cadres juridiques définissent la citation comme une reprise seconde et finalisée d'un extrait préexistant \cite{vivant_droit_2009,chambathouillon2016droit}. Les études médiatiques s'intéressent à la circulation de fragments discursifs autonomisés (slogans ou \emph{petites phrases}) dans l'espace public \cite{kriegplanque2010formule,boyer2018petites}.

Les approches linguistiques ne proposent pas de typologie directement
annotable de la citation, mais décrivent des phénomènes à un niveau
conceptuel plus général (opérations de reformulation, continuum des
formes, effets de point de vue). Ces descriptions constituent néanmoins
un socle théorique indispensable pour élaborer des typologies
opératoires adaptées aux corpus et aux méthodes computationnelles.

\section{Corpus existants pour l'étude et la détection des citations}

Les travaux en TAL ont donné lieu à de nombreux corpus annotés pour
l'étude de la citation et de l'attribution, reposant sur des conceptions
hétérogènes du discours rapporté.

Certains corpus, comme \textit{QuoteBank} (anglais), se concentrent
principalement sur les citations directes et l'identification du
locuteur, sans modéliser la diversité des formes de reprise.

À l'inverse, des schémas plus structurels, tels que PARC (anglais),
décrivent les relations d'attribution à partir de composantes explicites
(Source, Cue, Content) et distinguent citation directe, indirecte et
mixte a posteriori. Le corpus FRACAS (français) adopte une approche
similaire. Par exemple :

{[}\emph{Rihanna}{]}\textsubscript{speaker} {[}\emph{a
demandé}{]}\textsubscript{cue} {[}\emph{de ne pas arrêter la musique}{]}
\textsubscript{quote}.

\noindent La catégorie d'indirect y couvre toutefois des réalisations plus
hétérogènes.

PARC annote l'attribution de contenus relevant de la parole, de la
pensée ou de l'écrit, là où FRACAS regroupe ces réalisations sous des
catégories citationnelles plus larges.

Une ressource comme DE-News (allemand), adopte un schéma plus fin,
distinguant types de discours rapporté (dont le discours indirect libre)
et médias du dire (parole, pensée, écrit). Cette diversité reflète des
découpages théoriques distincts, que la typologie proposée vise à rendre
comparables.

\section{Pourquoi la définition de la citation devient un problème computationnel}

La définition de la citation n'est pas seulement un enjeu théorique :
elle conditionne le fonctionnement des systèmes de détection
automatique. Selon la famille de méthodes mobilisée, la notion de
citation est encodée, apprise ou inférée de manière différente, avec des
conséquences sur les résultats et leur interprétation.

Les systèmes à règles reposent sur l'identification de marqueurs formels
explicites, tels que les guillemets, la ponctuation ou des verbes
déclaratifs prototypiques, selon des patrons lexicosyntaxiques définis a
priori \cite{pouliquen2007automatic}. Ils détectent efficacement les
citations directes et certaines formes d'indirect, mais uniquement dans
le périmètre prévu par les règles : ce qui n'a pas été anticipé (formes
thématiques, indirect libre, reprises elliptiques) est mécaniquement
exclu. Ce type de pipeline reste largement utilisé dans des projets en
sciences humaines et sociales, notamment pour l'étude de la
représentation des genres dans la presse, comme dans \textit{GenderedNews}
\cite{richard2022genderednews} ou \textit{Radar de Parité} \cite{soumah2023radar}, où
l'extraction de citations constitue une étape intermédiaire du
traitement.

Les modèles discriminatifs, entraînés sur des corpus annotés, apprennent
à reconnaître les classes définies par ces corpus. Leur comportement
reflète la typologie locale de l'annotation, qu'il s'agisse de cadres
fondés sur l'attribution \cite{pareti2012database}, de corpus à grande échelle
comme \textit{QuoteBank} \cite{vaucher2021quotebank}, ou de ressources plus fines
telles que FRACAS \cite{richard2024fracas}. Plusieurs travaux comparatifs
montrent que les différences de performance entre approches à règles et
modèles neuronaux tiennent principalement aux choix de définition et
d'annotation de la citation, comme le montrent Janicki et al. \cite{janicki2023detection}
pour la presse finlandaise.

Les grands modèles de langage (LLMs) introduisent une situation
différente. Sans fine-tuning, ils mobilisent une notion implicite de la
citation issue de données d'entraînement hétérogènes ; après fine-tuning
sur un corpus donné, leur comportement se rapproche de celui des modèles
discriminatifs. La formulation du prompt encode alors implicitement une
définition de la citation, susceptible d'entrer en tension avec celle du
corpus.

Enfin, les scores sont comparables à l'intérieur d'un corpus, mais pas
entre ressources fondées sur des définitions différentes ; la typologie
proposée vise à rendre ces périmètres explicites.

\section{Notre proposition : vers une typologie unifiée et étendue des reprises discursives}

L'analyse de l'état de l'art et des corpus montre que la notion de
citation pose moins un problème terminologique qu'un problème de
périmètre. La typologie proposée vise à expliciter les formes de reprise
mobilisées dans les corpus, y compris en audiovisuel. Il s'agit d'un
schéma d'opérationnalisation. L'articulation générale de cette typologie
est synthétisée dans la Figure 1, qui présente l'organisation des
différentes formes de reprises discursives.

\begin{figure}[h]
\centering
\includegraphics[width=0.9\textwidth]{./media/image1.png}
\caption{Typologie des formes de reprise de discours.}
\end{figure}

Nous utilisons le terme de reprise de discours pour désigner un
sous-ensemble du discours rapporté correspondant aux formes donnant
accès à un dire ou à une pensée attribuable et annotable. Dans ce cadre,
nous excluons les discours narrativisés au sens narratologique \cite{genette1972figures}, qui se bornent à mentionner une activité discursive sans en
reprendre le contenu (\og il a expliqué sa position\fg{}), ainsi que les
formes de discours intérieur ou de pensée représentée qui ne permettent
pas, dans le contexte local, de reconstruire un dire ou une pensée
formulée \cite{cohn1978transparent}. Nous excluons également les formes
d'hétérogénéité énonciative \cite{authier1996ces}, qui signalent la
présence diffuse de discours autres sans attribution ni contenu
propositionnel identifiable (\og certains diront que\ldots{}\fg{}), ainsi que
les évaluations journalistiques du dire qui relèvent d'une
interprétation de posture discursive sans reprise effective (\og elle
dramatise la situation\fg{}).

Notre proposition se concentre sur la reprise d'un dire et repose sur un
principe simple : toutes les reprises ne relèvent pas du même régime
discursif. Nous distinguons deux domaines complémentaires : les
citations structurelles, fondées sur l'attribution explicite d'un dire,
et à des fins d'analyse médiatiques, les citations culturelles, fondées
sur la reconnaissance intertextuelle de fragments discursifs issus d'un
événement discursif identifié, circulant indépendamment de leur
attribution locale.

Les citations structurelles peuvent être continues ou discontinues
(séparées par des incises ou des commentaires narratifs), et porter sur
des propos issus de l'écrit, de l'oral ou de l'audiovisuel, repris dans
un texte, une transcription ou un document audiovisuel. Au sein de ce
domaine, nous proposons une typologie fondée sur la nature de la reprise
du dire : citation directe (reprise verbatim), citation indirecte
(reformulation d'un propos formulé), citation synthétique (reprise des
idées du propos) et citation mixte (combinaison de reprise littérale et
de reformulation).

Cette typologie est pleinement transversale aux supports. En
audiovisuel, la citation directe prend deux formes principales :
l'insert audio ou vidéo d'un propos antérieur, ou la reproduction
verbatim d'un propos par le narrateur. Dans le premier cas, seules les
reprises d'extraits déjà rendus publics constituent des reprises de
discours, ce qui exclut les prises de parole produites pour l'émission
(reportage original, micro-trottoir, direct) ou de leur rediffusion.
Dans le second cas, la citation est généralement signalée par des
marqueurs oraux (\og je cite\fg{}, \og selon ses propres mots\fg{}), ou par des
indices prosodiques ou multimodaux (variation du registre vocal, gestes
iconiques). La délimitation de la citation, notamment de sa fin,
constitue un enjeu saillant. En audiovisuel, elle repose sur des indices
combinés (retour à la voix initiale, rupture prosodique, changement
thématique, montage). À l'écrit, la clôture des citations indirectes
dépend souvent du réancrage énonciatif du narrateur plutôt que de
marqueurs formels.

Au sein des reprises structurellement marquées, nous distinguons enfin
une catégorie de propos interprété. Elle regroupe les cas où le
narrateur attribue à une source une position, une intention ou une
attitude discursive, sans rapporter explicitement un énoncé formulé,
mais en interprétant le dire ou la pensée de la source. Elle recouvre
notamment des attributions d'états mentaux (\og il pense que\ldots{}\fg{}),
des intentions discursives (\og elle veut convaincre\fg{}) et certaines
formes de cadrage journalistique du dire (\og il persiste et signe\fg{}),
lorsqu'elles fonctionnent comme attributions interprétatives, sans
ancrage explicite dans un dire rapporté. Cette distinction ne vise pas à
améliorer la détection, mais à isoler des cas interprétatifs difficiles
pour l'annotation et l'évaluation.

Les citations culturelles constituent des reprises non structurelles
relevant d'un régime de circulation intertextuelle (ex. « petites
phrases »), issues d'un dire historiquement situé et reconnues
indépendamment de leur attribution locale dans le document.

\section{Discussion: portée et usages d'une typologie opérationnelle}

La typologie proposée ne vise ni à remplacer les schémas d'annotation
existants ni à imposer un standard unique pour la détection de la
citation. Son principal apport est de rendre explicites et comparables
les choix théoriques structurant les corpus et les tâches
computationnelles. En distinguant la reprise d'un dire formulé (citation
structurelle) de l'interprétation narrative d'un dire (propos
interprété), elle permet de relire des ressources comme PARC, FRACAS ou
\textit{DE-News} non comme des ensembles incompatibles, mais comme des découpages
différents d'un même continuum du discours rapporté. Un test
d'alignement avec ces corpus permettrait ultérieurement d'en examiner la
couverture.

Cette explicitation est également déterminante pour l'interprétation des
résultats de détection automatique. Les divergences observées entre
systèmes à base de règles, modèles discriminatifs et modèles génératifs
apparaissent moins comme des limites algorithmiques que comme des effets
des définitions implicites de la citation encodées dans les données
d'entraînement et d'évaluation. Une typologie explicite permet de
distinguer désaccord théorique et difficulté technique, et d'envisager
des extensions de corpus cohérentes.

Un troisième enjeu concerne l'extension de la détection des citations à
l'audiovisuel. Les pratiques médiatiques contemporaines -- inserts de
discours, reprises verbatim orales, circulation de \og petites phrases\fg{}
médiatiques ou politiques -- ne peuvent être traitées sans une définition
explicite des objets annotés. L'introduction de la distinction entre
citations structurelles et citations culturelles constitue ici un apport
central : elle permet de séparer l'attribution locale d'un dire formulé
de la reconnaissance intertextuelle de fragments discursifs circulants
issus d'événements médiatiques. La définition proposée de la citation
directe audiovisuelle, par insert ou par reprise verbatim signalée,
fournit en outre une équivalence fonctionnelle aux guillemets de
l'écrit. Cette clarification ouvre la voie à la constitution de
nouvelles vérités terrain audiovisuelles et au développement de systèmes
de détection adaptés, tout en restant compatibles avec les corpus
textuels existants.

\section{Conclusion}\label{conclusion}

Cet article soutient que les difficultés de détection et d'évaluation
des citations tiennent moins aux méthodes qu'à l'instabilité
conceptuelle de l'objet \og citation\fg{}. Tant que cette notion reste
implicite ou variable selon les cadres théoriques, les performances des
systèmes demeurent difficiles à interpréter et à comparer.

Pour répondre à cette situation, nous avons proposé une typologie
opératoire qui articule plusieurs distinctions complémentaires. Au sein
des citations structurelles, la différenciation entre citations directe,
indirecte, synthétique et mixte permet de mieux circonscrire le
périmètre effectif de la reprise d'un dire. La séparation explicite
entre citation et propos interprété clarifie les frontières entre
attribution d'un propos formulé et interprétation narrative.
L'introduction des citations culturelles permet enfin de rendre compte
de la circulation médiatique de fragments discursifs emblématiques.

En intégrant ces distinctions et en prenant explicitement en compte
l'audiovisuel, notamment à travers les inserts et les reprises verbatim
orales, la typologie proposée fournit un cadre explicite pour comparer
les choix d'annotation, guider l'extension des corpus et articuler
définition linguistique, détection automatique et évaluation.

%\nocite{*}
\printbibliography

\end{document}